\clearpage
\begin{trapbox}
\textbf{C090\quad 易错清单、边界与最小反例}

\begin{center}
% At P=(2,3), ellipse tangent slope=-1/2; hyperbola tangent slope=2.
\begin{tikzpicture}[x=.9cm,y=.9cm,every node/.style={font=\small}]
\path[use as bounding box] (-3,-.8) rectangle (8,5.2);
\coordinate (P) at (2,3); \coordinate (L) at (-2,0); \coordinate (R) at (2,0);
\draw[teal!70!black,very thick] (L)--(P)--(R);
\draw[gray] (L)--(R);
\draw[blue!70!black,thick,dashed] (-.5,4.25)--(4.5,1.75);
\draw[red!70!black,thick,dash dot] (.5,0)--(3,5);
\draw[teal!70!black] (2,2.4) arc(-90:-143.130102:.6);
\node at (1.05,2) {$\theta$};
\draw (2.2683,2.8658)--(2.4025,3.1342)--(2.1342,3.2683);
\fill (P) circle(2pt);
\node at (3.4,3.7) {$P(2,3)$};
\fill (L) circle(1.6pt) node[below] {$F_1$};
\fill (R) circle(1.6pt) node[below] {$F_2$};
\node[blue!70!black,above] at (-.1,4.05) {椭圆切线};
\node[red!70!black,right] at (3,4.8) {双曲线切线};
\node[anchor=west,align=left] at (4.8,2.5) {焦半径夹角：\\[5pt] $\cos\theta=3/5$\\[12pt] 切线夹角：$90^\circ$};
\end{tikzpicture}

\end{center}
\begin{itemize}
\item \textbf{混淆两曲线的半长轴。}椭圆 $a_E=c/e_1>c$，双曲线 $a_H=c/e_2<c$。
不能把两处的 $a$ 当成同一个参数。最稳妥是记作 $A$ 与 $T$。

\item \textbf{焦半径差忘记绝对值。}双曲线定义是
$|PF_1-PF_2|=2a_H$，左支与右支交换正负号；主公式只含差的平方，
因此不受支别影响。

\item \textbf{余弦项符号颠倒。}正确的是
\[
 \boxed{(1-\cos\theta)/e_1^2+(1+\cos\theta)/e_2^2=2}.
\]
记忆：\emph{和}对应 $1-\cos\theta$（椭圆），\emph{差}对应 $1+\cos\theta$（双曲线）。

\item \textbf{把两种夹角混为一谈。}$\theta$ 是 $\angle F_1PF_2$，
图中实线为焦半径，虚线与点划线为两条切线；后者正交不代表前者正交。

\item \textbf{误判柯西等号条件。}向量配比要求
$\sin^2(\theta/2)/e_1=\cos^2(\theta/2)/e_2$，
不是 $e_1=e_2$；由于一小于 $1$ 一大于 $1$，$e_1=e_2$ 根本不可能。

\item \textbf{误以为任意角度都能取到推论的等号。}等号对应
$e_1=\tan(\theta/2)<1$、$e_2=\cot(\theta/2)>1$，
要求 $\theta<90^\circ$；$\theta=120^\circ$ 时不等式严格。

\item \textbf{误用 $\theta=0$ 或 $\pi$。}公共点不在两焦点所在直线上，
所以 $0<\theta<\pi$ 且 $\sin\theta>0$。推论不能允许端点角度。

\item \textbf{对固定角度漏掉可行边界。}必须
$\sin(\theta/2)<e_1<1$；例如 $\theta=60^\circ$ 时
$e_1=1/2$ 不可行；此时第一项已等于 $1$，却还有正的第二项。

\item \textbf{误把四个交点认成一个。}对任意 $0<e_1<1<e_2$ 的标准共焦
椭圆、双曲线，公共点有四个，坐标分别为 $(\pm X,\pm Y)$（$X,Y>0$）；
四点的焦角相同，正交对称性不能被误作只在右支成立。
\end{itemize}
\end{trapbox}

\clearpage
\begin{trapbox}
\textbf{从学生最常见的疑问出发：四组最小反例与边界。}
\begin{itemize}
\item \textbf{错误：看到示意图有直角，就假设所有共焦交点都形成直角三角形。}
反例：取 $A=5,T=2,c=3$，焦点三角形三边是 $7,3,6$，
$\cos\theta=11/21\ne0$。\emph{一般证明须使用一般三角形}；
$3,4,5$ 直角图只用于演示特殊的等号情况。

\item \textbf{错误：把 $e_1e_2=1$ 当作“$\theta=90^\circ$”。}
正确的是：在共焦、有公共点前提下，$e_1e_2=1$
等价于焦角不等式取到等号；而\emph{直角位于一个焦点}，
$\theta=\angle F_1PF_2$ 实际是\emph{锐角}。
以焦点为 $(\pm c,0)$，由 $x^2=A^2T^2/c^2$ 可见
$AT=c^2$ 时 $x=\pm c$，公共点就在过某个焦点的垂线上。

\item \textbf{错误：$\theta=120^\circ$ 时也认为
$2/\sin120^\circ$ 是可取到的最大值。}
这是\emph{有效但不紧的上界}。当 $\theta\ge90^\circ$，
实际上 $1/e_1+1/e_2<2$，可无限接近 $2$ 而不能取到；
详见例6。\emph{“有上界”不等于“存在最大值”。}

\item \textbf{错误：用曲线切线的 $90^\circ$ 代替焦角 $\theta$。}
共焦椭圆与双曲线在公共点处确实正交，
但它们的切线夹角与 $\angle F_1PF_2$ 是不同的角。
对于例1，切线夹角是 $90^\circ$，焦角却满足
$\cos\theta=3/5$。
\end{itemize}
\end{trapbox}
